\documentclass[twocolumn,showpacs,preprintnumbers,amsmath,amssymb,prl]{revtex4-2}

\usepackage{amsmath}
\usepackage{amssymb}
\usepackage{amsthm}
\usepackage{graphicx}
\usepackage{hyperref}
% physics package replaced with manual macros (not in BasicTeX by default)
\newcommand{\ket}[1]{|#1\rangle}
\newcommand{\bra}[1]{\langle #1|}
\newcommand{\braket}[2]{\langle #1|#2\rangle}
\newcommand{\tr}{\mathrm{Tr}}
\newcommand{\abs}[1]{\left|#1\right|}

\newtheorem{lemma}{Lemma}

% ------------------------------------------------------------------
% Convenience macros
% ------------------------------------------------------------------
\newcommand{\K}{\mathcal{K}}          % KCBS witness
\newcommand{\pstar}{p^{*}}            % noise threshold
\newcommand{\rhobar}{\bar{\rho}}      % noisy state
\newcommand{\Id}{\mathbb{I}}          % identity

\begin{document}

\title{KCBS Contextuality Under Non-Depolarizing Noise:\\
       Analytic Thresholds for Dephasing and Amplitude Damping}

\author{Mario Ivanov}
\email{mariohgh@yahoo.co.uk}
\affiliation{Independent researcher. ORCID: 0009-0005-3167-2162}

\date{\today}

% ==================================================================
\begin{abstract}
We derive closed-form expressions for the Klyachko--Can--Binicio\u{g}lu--Shumovsky
(KCBS) witness $\K(p)$ under two noise channels that have not previously been
analyzed analytically for a spin-1 qutrit: dephasing and full-cascade amplitude
damping.
For the KCBS-optimal state $\ket{2}\bra{2}$, dephasing leaves the witness
invariant ($\K(p) = \sqrt{5}$ for all $p$) because the optimal state carries no
coherences.
Amplitude damping yields $\K(p) = A\,p(2-p) + \sqrt{5}\,(1-p)^2$, where
$A = \tfrac{5}{2} - \tfrac{\sqrt{5}}{2}$, with noise threshold
$\pstar_{\mathrm{AD}} = 1 - \sqrt{(\sqrt{5}-1)/(3\sqrt{5}-5)} \approx 0.149$.
This is $2.78\times$ lower than the depolarizing threshold
$\pstar_{\mathrm{dep}} \approx 0.415$, showing that the geometry of the noise
channel --- not only its strength --- determines how quickly contextuality is lost.
All results are verified numerically to machine precision.
\end{abstract}

\maketitle

% ==================================================================
\section{Introduction}
\label{sec:intro}

Quantum contextuality --- the impossibility of assigning predetermined outcomes
to measurements independently of which compatible measurements are jointly
performed --- is a foundational non-classical feature of quantum
mechanics~\cite{KochenSpecker1967,Budroni2022}.
The Klyachko--Can--Binicio\u{g}lu--Shumovsky (KCBS) inequality provides a
particularly clean operational test for a spin-1 (qutrit) system: five
sequential measurements whose correlations satisfy $\K \leq 2$ for any
noncontextual model, while quantum mechanics predicts $\K = \sqrt{5} \approx
2.236$ for an optimal pure state~\cite{Klyachko2008}.

Real experiments are subject to noise, and a natural question is: how much noise
can the system tolerate before the quantum violation disappears?
For depolarizing noise, the answer is known analytically~\cite{Klyachko2008,Dutra2026}:
the KCBS witness decreases linearly with the noise parameter and the violation
survives up to $\pstar_{\mathrm{dep}} \approx 0.415$.
Shaham and Eisenberg tested dephasing and two other unital channels experimentally
on a biphoton qutrit, finding that dephasing suppresses the violation at a noise
level of approximately $P \sim 0.17$ (read from a graph; no analytic formula was
given)~\cite{Shaham2014}.
Dutra, Baldij\~{a}o, and Terra Cunha recently extended the depolarizing analysis
to sequential measurement noise in both KCBS and Peres--Mermin scenarios, and
explicitly named dephasing as a physically relevant alternative channel without
analyzing it~\cite{Dutra2026}.

In this paper we fill that gap.
We derive analytic expressions for $\K(p)$ under dephasing and full-cascade
amplitude damping on a single spin-1 qutrit, identify the noise thresholds in
closed form, and compare all three channels.
The results show that the structure of the noise --- specifically, whether it
acts on coherences, populations, or both --- determines the robustness of
contextuality in a geometrically transparent way.

% ==================================================================
\section{Setup}
\label{sec:setup}

\subsection{KCBS scenario}

The system is a qutrit with Hilbert space $\mathcal{H} = \mathbb{C}^3$.
Following Klyachko et al.~\cite{Klyachko2008}, we define five rank-one projectors
$P_i = \ket{v_i}\bra{v_i}$, $i = 0,\ldots,4$, whose unit vectors lie on a cone:
\begin{equation}
  \ket{v_i} = \sin\theta\cos\phi_i\,\ket{0}
             + \sin\theta\sin\phi_i\,\ket{1}
             + \cos\theta\,\ket{2},
  \label{eq:projectors}
\end{equation}
where $\cos^2\theta = 1/\sqrt{5}$ and $\phi_i = 4\pi i/5$.
Adjacent projectors satisfy $P_i P_{i+1\,\mathrm{mod}\,5} = 0$.

The KCBS witness is
\begin{equation}
  \K(\rho) = \sum_{i=0}^{4} \tr(\rho\, P_i).
  \label{eq:witness}
\end{equation}
Any noncontextual hidden-variable model satisfies $\K \leq 2$.
The optimal quantum state is $\rho_0 = \ket{2}\bra{2}$, which gives
$\K(\rho_0) = \sqrt{5}$.

\subsection{Noise model}

We model noise as a completely positive trace-preserving (CPTP) map acting on
the state before measurement.
This is physically distinct from noise acting between sequential measurements;
we discuss this distinction in Sec.~\ref{sec:discussion}.
Three channels are compared:

\paragraph{Depolarizing.}
\begin{equation}
  \mathcal{D}_p(\rho) = (1-p)\,\rho + p\,\frac{\Id}{3}, \quad p \in [0,1].
  \label{eq:depolarize}
\end{equation}

\paragraph{Dephasing.}
\begin{equation}
  \mathcal{Z}_p(\rho)_{ij} =
  \begin{cases}
    \rho_{ij} & i = j, \\
    (1-p)\,\rho_{ij} & i \neq j.
  \end{cases}
  \label{eq:dephase}
\end{equation}
This is a unital channel that preserves populations and suppresses coherences.

\paragraph{Full-cascade amplitude damping.}
Models energy decay $\ket{2}\to\ket{1}\to\ket{0}$ with equal decay probability
$p$ per level.
The Kraus operators are
\begin{align}
  K_0 &= \mathrm{diag}(1,\,\sqrt{1-p},\,(1-p)), \notag \\
  K_1 &= \sqrt{p}\,\ket{0}\bra{1}, \notag \\
  K_2 &= \sqrt{2p(1-p)}\,\ket{1}\bra{2}, \notag \\
  K_3 &= p\,\ket{0}\bra{2}.
  \label{eq:kraus}
\end{align}
Completeness $\sum_k K_k^\dagger K_k = \Id$ follows from the binomial identity
$(1-p)^2 + 2p(1-p) + p^2 = 1$.

% ==================================================================
\section{Results}
\label{sec:results}

\subsection{Depolarizing channel (baseline)}
\label{sec:depolarizing}

Substituting Eq.~\eqref{eq:depolarize} into Eq.~\eqref{eq:witness} and using
$\tr(\Id/3 \cdot P_i) = 1/3$ for any rank-one projector gives
\begin{equation}
  \K_{\mathrm{dep}}(p) = (1-p)\sqrt{5} + p\,\frac{5}{3}.
  \label{eq:K_dep}
\end{equation}
Setting $\K_{\mathrm{dep}}(\pstar) = 2$ yields
\begin{equation}
  \pstar_{\mathrm{dep}} = \frac{\sqrt{5}-2}{\sqrt{5}-5/3} \approx 0.4146.
  \label{eq:pstar_dep}
\end{equation}
This reproduces the known result~\cite{Klyachko2008,Dutra2026}.

\subsection{Dephasing channel}
\label{sec:dephasing}

\begin{lemma}
  $\K_{\mathrm{deph}}(p) = \sqrt{5}$ for all $p \in [0,1]$.
  The KCBS violation is fully robust to dephasing on the optimal state.
\end{lemma}

\begin{proof}
  The optimal state $\rho_0 = \ket{2}\bra{2}$ is diagonal in the standard
  basis: $(\rho_0)_{ij} = \delta_{i2}\delta_{j2}$.
  All off-diagonal elements are zero, so $\mathcal{Z}_p(\rho_0) = \rho_0$ for
  all $p$.
  Therefore $\K(\mathcal{Z}_p(\rho_0)) = \K(\rho_0) = \sqrt{5}$.
\end{proof}

This result is specific to the KCBS-optimal state.
A state with nonzero coherences would be affected by dephasing; the immunity
here reflects the geometric fact that $\ket{2}$ is the north pole of the
projector cone defined in Eq.~\eqref{eq:projectors}.

\subsection{Full-cascade amplitude damping}
\label{sec:amplitude_damping}

Applying the channel~\eqref{eq:kraus} to $\rho_0 = \ket{2}\bra{2}$ gives a
diagonal output state (amplitude damping preserves diagonality when the input
is diagonal):
\begin{equation}
  \mathcal{A}_p(\rho_0) = p^2\ket{0}\bra{0}
                        + 2p(1-p)\ket{1}\bra{1}
                        + (1-p)^2\ket{2}\bra{2}.
  \label{eq:rho_ad}
\end{equation}

To evaluate the witness, we need $\sum_i \abs{v_i^{(k)}}^2$ for each component
$k \in \{0,1,2\}$ of the projector vectors.
From Eq.~\eqref{eq:projectors}:
\begin{align}
  \sum_{i=0}^{4} \abs{v_i^{(0)}}^2
    &= \sum_{i=0}^{4} \abs{v_i^{(1)}}^2
     = \sin^2\theta \cdot \frac{5}{2}
     = \frac{5}{2} - \frac{\sqrt{5}}{2} \equiv A, \label{eq:sumA} \\
  \sum_{i=0}^{4} \abs{v_i^{(2)}}^2
    &= 5\cos^2\theta = \sqrt{5}. \label{eq:sumB}
\end{align}
The azimuthal sums $\sum_i \cos^2\phi_i = \sum_i \sin^2\phi_i = 5/2$ follow
from the five-fold symmetry of the $\phi_i = 4\pi i/5$ angles.

Substituting Eq.~\eqref{eq:rho_ad} into Eq.~\eqref{eq:witness}:
\begin{align}
  \K_{\mathrm{AD}}(p)
    &= p^2 \cdot A + 2p(1-p) \cdot A + (1-p)^2 \cdot \sqrt{5} \notag \\
    &= A\,p(2-p) + \sqrt{5}\,(1-p)^2,
  \label{eq:K_ad}
\end{align}
where $A = 5/2 - \sqrt{5}/2$.

Setting $\K_{\mathrm{AD}}(\pstar) = 2$ gives the quadratic
$(B-A)\,p^2 - 2(B-A)\,p + (B-2) = 0$
with $B = \sqrt{5}$, $C \equiv B - A = 3\sqrt{5}/2 - 5/2$.
The physical root (in $[0,1]$) is
\begin{equation}
  \pstar_{\mathrm{AD}}
    = 1 - \sqrt{\frac{\sqrt{5}-1}{3\sqrt{5}-5}}
    \approx 0.1493.
  \label{eq:pstar_ad}
\end{equation}

% ==================================================================
\section{Comparison and Discussion}
\label{sec:discussion}

\begin{figure}[h]
  \centering
  \includegraphics[width=\columnwidth]{../kcbs_three_channels.png}
  \caption{KCBS witness $\K(p)$ under three noise channels on the optimal
           qutrit state $\ket{2}\bra{2}$.
           Shaded regions indicate quantum violation ($\K > 2$).
           Dashed vertical lines mark the noise thresholds for depolarizing
           ($\pstar \approx 0.415$) and amplitude damping ($\pstar \approx 0.149$).}
  \label{fig:channels}
\end{figure}

The three results are summarized in Table~\ref{tab:thresholds} and
Fig.~\ref{fig:channels}.

\begin{table}[h]
  \centering
  \begin{tabular}{lcc}
    \hline\hline
    Channel & $\K(p)$ & $\pstar$ \\
    \hline
    Depolarizing     & $(1-p)\sqrt{5} + p\cdot\tfrac{5}{3}$          & $0.4146$ \\
    Dephasing        & $\sqrt{5}$ (constant)                          & none     \\
    Amplitude damping & $A\,p(2-p)+\sqrt{5}(1-p)^2$                  & $0.1493$ \\
    \hline\hline
  \end{tabular}
  \caption{KCBS witness and noise threshold for each channel on $\rho_0 = \ket{2}\bra{2}$.
           $A = 5/2 - \sqrt{5}/2 \approx 1.382$.}
  \label{tab:thresholds}
\end{table}

The ratio $\pstar_{\mathrm{dep}}/\pstar_{\mathrm{AD}} \approx 2.78$ shows that
amplitude damping destroys contextuality nearly three times faster than
depolarization at equal noise parameter values.
The physical reason is transparent from the geometry of the KCBS scenario.
The witness $\K$ is maximized by $\ket{2}\bra{2}$, which sits at the north pole
of the projector cone.
Depolarizing noise moves the state uniformly toward the maximally mixed state
$\Id/3$, which still has $\K(\Id/3) = 5/3 > 0$ --- a relatively mild
degradation.
Amplitude damping funnels all population toward the ground state $\ket{0}\bra{0}$,
which is the antipodal point on the cone and gives $\K(\ket{0}\bra{0}) = A \approx
1.382$ --- well below the noncontextual bound.
The state therefore crosses the bound much sooner.

Dephasing is inert on $\rho_0$ for a different reason: it acts only on
off-diagonal elements, and $\rho_0$ has none.
This is not a general statement about dephasing and KCBS contextuality --- a
superposition state would be affected --- but it is a precise statement about
the optimal state, and it has a direct experimental implication: an experiment
that prepares $\ket{2}$ and applies dephasing noise will not see any degradation
of the KCBS witness due to that noise alone.

\paragraph{Noise model caveat.}
All results here model noise as acting on the quantum state before measurement
(Schr\"{o}dinger picture, state noise).
This is physically distinct from noise acting between sequential measurements,
which was analyzed for the depolarizing channel in~\cite{Dutra2026}.
The two situations can give different thresholds and represent different
experimental scenarios; the distinction must be stated explicitly in any
experimental implementation.

\paragraph{Comparison with Shaham and Eisenberg~\cite{Shaham2014}.}
Their experiment found that dephasing suppresses KCBS contextuality at
$P \sim 0.17$ on a biphoton qutrit.
Our result --- that dephasing has no effect on the spin-1 optimal state ---
is not in contradiction.
The biphoton system is physically different: noise acts independently on each
photon, not on the qutrit as a whole, and the KCBS projectors are constructed
differently.
The two results together show that the robustness of contextuality to dephasing
depends sensitively on the physical realization of the qutrit.

% ==================================================================
\section{Conclusion}
\label{sec:conclusion}

We have derived analytic expressions for the KCBS witness under dephasing and
full-cascade amplitude damping on a spin-1 qutrit.
The dephasing result is a lemma: the KCBS-optimal state is immune to dephasing
because it carries no coherences.
The amplitude damping result is the main contribution: a closed-form $\K(p)$
and threshold $\pstar_{\mathrm{AD}} \approx 0.149$, derived from the binomial
cascade structure of the Kraus operators and the five-fold symmetry of the KCBS
projectors.

The comparison across three channels shows that the geometry of the noise ---
which degrees of freedom it acts on --- matters as much as its strength.
Amplitude damping is nearly three times more destructive to KCBS contextuality
than depolarization at equal noise parameter, while dephasing is completely
harmless on the optimal state.

These results provide analytic benchmarks for experimental implementations of
KCBS tests in systems where amplitude damping is the dominant noise process,
such as nitrogen-vacancy centers, trapped ions, and superconducting transmons
operating as effective qutrits.

\begin{acknowledgments}
Computations were performed using NumPy~\cite{NumPy2020}.
Code and data are available at \url{https://github.com/mariohgh/kcbs-noise-thresholds}.
This preprint is archived at \url{https://doi.org/10.5281/zenodo.23415822}.
\end{acknowledgments}

\bibliography{refs}

\end{document}
